% template-assignment.tex --- Example assignment using academic-report class
\documentclass[assignment]{academic-report}

% --- Metadata ---
\title{Problem Set 3: Entanglement Measures}
\author{Student Name}
\date{\today}

% --- Course Info ---
\course{PHYS 641: Advanced Quantum Mechanics}
\instructor{Prof. S. Chen}
\duedate{October 15, 2026}

\begin{document}

\maketitle

\section{Entanglement Entropy of Bipartite Systems}

\begin{problem}[Entanglement entropy of a two-spin system]
Consider two spin-$\frac{1}{2}$ particles in the singlet state:
\begin{equation}
  |\psi\rangle = \frac{1}{\sqrt{2}}\left(|\uparrow\downarrow\rangle -
  |\downarrow\uparrow\rangle\right)
\end{equation}
Compute the entanglement entropy of one spin by tracing out the other.
\end{problem}

\begin{solution}
The density matrix of the full system is:
\begin{equation}
  \rho = |\psi\rangle\langle\psi| =
  \frac{1}{2}\left(|\uparrow\downarrow\rangle\langle\uparrow\downarrow| +
  |\downarrow\uparrow\rangle\langle\downarrow\uparrow| -
  |\uparrow\downarrow\rangle\langle\downarrow\uparrow| -
  |\downarrow\uparrow\rangle\langle\uparrow\downarrow|\right)
\end{equation}

Tracing out the second spin yields the reduced density matrix:
\begin{equation}
  \rho_A = \mathrm{Tr}_B(\rho) =
  \frac{1}{2}\left(|\uparrow\rangle\langle\uparrow| +
  |\downarrow\rangle\langle\downarrow|\right) =
  \begin{pmatrix} \frac{1}{2} & 0 \\ 0 & \frac{1}{2} \end{pmatrix}
\end{equation}

The entanglement entropy is:
\begin{equation}
  S_A = -\mathrm{Tr}(\rho_A \ln \rho_A) =
  -2 \cdot \frac{1}{2}\ln\frac{1}{2} = \ln 2
\end{equation}

This is the maximum possible entanglement for a two-qubit system.
\end{solution}

\begin{problem}[Area law in the AKLT chain]
The Affleck-Kennedy-Lieb-Tasaki (AKLT) chain has the following Hamiltonian:
\begin{equation}
  H = \sum_i \left[\mathbf{S}_i \cdot \mathbf{S}_{i+1} +
  \frac{1}{3}(\mathbf{S}_i \cdot \mathbf{S}_{i+1})^2\right]
\end{equation}

Using the matrix product state representation, argue that the ground state
entanglement entropy saturates to a constant $S = 2\ln 2$ for large block
sizes, consistent with the area law in one dimension.
\end{problem}

\begin{solution}
The AKLT ground state can be written as an MPS with bond dimension $\chi = 2$.
The virtual space at each bond encodes a spin-$\frac{1}{2}$ degree of freedom,
so the Schmidt decomposition at any cut has exactly two non-zero singular
values $\lambda_1 = \lambda_2 = 1/\sqrt{2}$.

For a block of length $L$, the entanglement entropy is bounded by:
\begin{equation}
  S_A(L) \leq \ln(\chi^2) = \ln 4 = 2\ln 2
\end{equation}

For large $L$, this bound is saturated. Since the entropy is independent of
block length, it satisfies the one-dimensional area law: $S_A \sim \text{const}$
for gapped systems.
\end{solution}

\section{Computational Methods}

\begin{problem}[Numerical entanglement spectrum]
Write a short program to compute the entanglement spectrum of the transverse
field Ising model at criticality. The Hamiltonian is:
\begin{equation}
  H = -J\sum_i \left(\sigma_i^z \sigma_{i+1}^z + g\sigma_i^x\right)
\end{equation}
with $g = 1$ at the critical point.
\end{problem}

\begin{solution}
The following code uses exact diagonalization for a small chain:

\begin{lstlisting}[language=Python, caption={Exact diagonalization of TFI model}]
import numpy as np
from scipy.linalg import eigh

def tfi_hamiltonian(N, J=1.0, g=1.0):
    """Build the transverse field Ising Hamiltonian."""
    dim = 2**N
    H = np.zeros((dim, dim))

    for i in range(N):
        j = (i + 1) % N
        # sigma^z_i sigma^z_{j+1}
        for s in range(dim):
            spin_i = (s >> i) & 1
            spin_j = (s >> j) & 1
            H[s, s] += -J * (1 - 2*spin_i) * (1 - 2*spin_j)
        # g * sigma^x_i
        for s in range(dim):
            s_flip = s ^ (1 << i)
            H[s, s_flip] += -J * g

    return H

N = 10
H = tfi_hamiltonian(N, g=1.0)
energies, states = eigh(H)

# Ground state
psi_0 = states[:, 0]
print(f"Ground state energy: {energies[0]:.6f}")
print(f"Energy gap: {energies[1] - energies[0]:.6f}")
\end{lstlisting}

The entanglement spectrum can then be obtained by reshaping the ground state
vector into a matrix and performing SVD.
\end{solution}

\end{document}
